\documentclass[10pt,a4paper]{amsart}
\usepackage{amsmath,amssymb,amsthm,graphicx}
% Retained minimal preamble for the Stacks Project chapter more-morphisms.tex (More on Morphisms).
% Verbatim excerpt of the chapter preamble: Xy-pic setup (native preamble lines 14-18) and the native
% theorem declarations with the mathematical macros (native preamble lines 154-217). Copied mechanically;
% nothing added, removed or redefined. The wrapper class and the amsmath/amssymb/amsthm/graphicx base
% packages are supplied by the assembler, so the class/externaldocument part of the native preamble is omitted.
% For commutative diagrams we use Xy-pic
\usepackage[all]{xy}

% We use 2cell for 2-commutative diagrams.
\xyoption{2cell}
% Theorem environments.
%
\theoremstyle{plain}
\newtheorem{theorem}[subsection]{Theorem}
\newtheorem{proposition}[subsection]{Proposition}
\newtheorem{lemma}[subsection]{Lemma}

\theoremstyle{definition}
\newtheorem{definition}[subsection]{Definition}
\newtheorem{example}[subsection]{Example}
\newtheorem{exercise}[subsection]{Exercise}
\newtheorem{situation}[subsection]{Situation}

\theoremstyle{remark}
\newtheorem{remark}[subsection]{Remark}
\newtheorem{remarks}[subsection]{Remarks}

\numberwithin{equation}{subsection}

% Macros
%
\def\lim{\mathop{\mathrm{lim}}\nolimits}
\def\colim{\mathop{\mathrm{colim}}\nolimits}
\def\Spec{\mathop{\mathrm{Spec}}}
\def\Hom{\mathop{\mathrm{Hom}}\nolimits}
\def\Ext{\mathop{\mathrm{Ext}}\nolimits}
\def\SheafHom{\mathop{\mathcal{H}\!\mathit{om}}\nolimits}
\def\SheafExt{\mathop{\mathcal{E}\!\mathit{xt}}\nolimits}
\def\Sch{\mathit{Sch}}
\def\Mor{\mathop{\mathrm{Mor}}\nolimits}
\def\Ob{\mathop{\mathrm{Ob}}\nolimits}
\def\Sh{\mathop{\mathit{Sh}}\nolimits}
\def\NL{\mathop{N\!L}\nolimits}
\def\CH{\mathop{\mathrm{CH}}\nolimits}
\def\proetale{{pro\text{-}\acute{e}tale}}
\def\etale{{\acute{e}tale}}
\def\QCoh{\mathit{QCoh}}
\def\Ker{\mathop{\mathrm{Ker}}}
\def\Im{\mathop{\mathrm{Im}}}
\def\Coker{\mathop{\mathrm{Coker}}}
\def\Coim{\mathop{\mathrm{Coim}}}

% Boxtimes
%
\DeclareMathSymbol{\boxtimes}{\mathbin}{AMSa}{"02}

%
% Macros for moduli stacks/spaces
%
\def\QCohstack{\mathcal{QC}\!\mathit{oh}}
\def\Cohstack{\mathcal{C}\!\mathit{oh}}
\def\Spacesstack{\mathcal{S}\!\mathit{paces}}
\def\Quotfunctor{\mathrm{Quot}}
\def\Hilbfunctor{\mathrm{Hilb}}
\def\Curvesstack{\mathcal{C}\!\mathit{urves}}
\def\Polarizedstack{\mathcal{P}\!\mathit{olarized}}
\def\Complexesstack{\mathcal{C}\!\mathit{omplexes}}
% \Pic is the operator that assigns to X its picard group, usage \Pic(X)
% \Picardstack_{X/B} denotes the Picard stack of X over B
% \Picardfunctor_{X/B} denotes the Picard functor of X over B
\def\Pic{\mathop{\mathrm{Pic}}\nolimits}
\def\Picardstack{\mathcal{P}\!\mathit{ic}}
\def\Picardfunctor{\mathrm{Pic}}
\def\Deformationcategory{\mathcal{D}\!\mathit{ef}}
\newenvironment{reference}{}{}
\newenvironment{slogan}{}{}
\newenvironment{history}{}{}
\usepackage{cleveref}
\crefname{originallemma}{original lemma}{original lemmas}
\Crefname{originallemma}{Original lemma}{Original lemmas}
\makeatletter
\newlabel{algebra-lemma-etale-makes-quasi-finite-finite-variant}{{algebra-lemma-etale-makes-quasi-finite-finite-variant}{source}{}{sourcepointer-1}{}}
\newlabel{algebra-lemma-etale-makes-quasi-finite-finite-variant@cref}{{[originallemma][0][]algebra-lemma-etale-makes-quasi-finite-finite-variant}{[1][0][]source}{}{}{}}
\makeatother
\crefname{originallemma}{original lemma}{original lemmas}
\Crefname{originallemma}{Original lemma}{Original lemmas}
\makeatletter
\newlabel{lemma-etale-makes-quasi-finite-finite-multiple-points}{{lemma-etale-makes-quasi-finite-finite-multiple-points}{source}{}{sourcepointer-2}{}}
\newlabel{lemma-etale-makes-quasi-finite-finite-multiple-points@cref}{{[originallemma][0][]lemma-etale-makes-quasi-finite-finite-multiple-points}{[1][0][]source}{}{}{}}
\makeatother
\crefname{originallemma}{original lemma}{original lemmas}
\Crefname{originallemma}{Original lemma}{Original lemmas}
\makeatletter
\newlabel{lemma-realize-prescribed-residue-field-extension-etale}{{lemma-realize-prescribed-residue-field-extension-etale}{source}{}{sourcepointer-3}{}}
\newlabel{lemma-realize-prescribed-residue-field-extension-etale@cref}{{[originallemma][0][]lemma-realize-prescribed-residue-field-extension-etale}{[1][0][]source}{}{}{}}
\makeatother
\crefname{originallemma}{original lemma}{original lemmas}
\Crefname{originallemma}{Original lemma}{Original lemmas}
\makeatletter
\newlabel{morphisms-lemma-base-change-quasi-finite}{{01TM}{source}{}{sourcepointer-4}{}}
\newlabel{morphisms-lemma-base-change-quasi-finite@cref}{{[originallemma][0][]01TM}{[1][0][]source}{}{}{}}
\makeatother
\crefname{originallemma}{original lemma}{original lemmas}
\Crefname{originallemma}{Original lemma}{Original lemmas}
\makeatletter
\newlabel{morphisms-lemma-quasi-finite-at-point-characterize}{{01TH}{source}{}{sourcepointer-5}{}}
\newlabel{morphisms-lemma-quasi-finite-at-point-characterize@cref}{{[originallemma][0][]01TH}{[1][0][]source}{}{}{}}
\makeatother
\crefname{originallemma}{original lemma}{original lemmas}
\Crefname{originallemma}{Original lemma}{Original lemmas}
\makeatletter
\newlabel{morphisms-lemma-residue-field-quasi-finite}{{morphisms-lemma-residue-field-quasi-finite}{source}{}{sourcepointer-6}{}}
\newlabel{morphisms-lemma-residue-field-quasi-finite@cref}{{[originallemma][0][]morphisms-lemma-residue-field-quasi-finite}{[1][0][]source}{}{}{}}
\makeatother
\crefname{originallemma}{original lemma}{original lemmas}
\Crefname{originallemma}{Original lemma}{Original lemmas}
\makeatletter
\newlabel{schemes-lemma-points-fibre-product}{{schemes-lemma-points-fibre-product}{source}{}{sourcepointer-7}{}}
\newlabel{schemes-lemma-points-fibre-product@cref}{{[originallemma][0][]schemes-lemma-points-fibre-product}{[1][0][]source}{}{}{}}
\makeatother
\title{Making several isolated fibre points finite by an etale neighbourhood of several points}
\author{The Stacks Project Authors}
\date{stacks-project:more-morphisms.tex:local-snapshot-20261011; GNU FDL 1.2 or later}
\begin{document}
\maketitle
\noindent\textit{Editorial note.} Statement span is lines 12018-12044 of sources/stacks-more-morphisms/pinned-original-source.tex; proof span is lines 12046-12091, the complete author proof (44 body lines).\par\medskip
\noindent\textit{Editorial note.} Navigation stubs: 7 cited-but-not-included labels are supplied as non-load-bearing pointers with their pinned tags where the snapshot tag index has them.\par\medskip
\noindent\textit{Editorial note.} extra\_packages is empty on purpose: the retained preamble already loads xy with options and a second declaration clashes in the assembler wrapper.\par\medskip
\noindent\textit{Editorial note.} Text is retained verbatim, including the author's own typos and placeholders.\par\medskip
\noindent\textit{Allowed standard prerequisites (editorial).}\begin{itemize}
\item Explicit prerequisite boundary. Etale neighbourhoods, isolated points in a fibre, and finiteness of a morphism over an etale base change; the assembly is retained.
\end{itemize}\medskip
\ifdefined\hypertarget\hypertarget{sourcepointer-1}{}\fi
\noindent\textit{Source reference.} Original lemma algebra-lemma-etale-makes-quasi-finite-finite-variant: sources/stacks-more-morphisms/pinned-original-source.tex, source lines 12091--12091. Cited by the retained statement or proof but not part of the collected spans; supplied only as a navigation pointer. The argument does not depend on its content.\par\medskip
\ifdefined\hypertarget\hypertarget{sourcepointer-2}{}\fi
\noindent\textit{Source reference.} Original lemma lemma-etale-makes-quasi-finite-finite-multiple-points: sources/stacks-more-morphisms/pinned-original-source.tex, source lines 11974--11981. Cited by the retained statement or proof but not part of the collected spans; supplied only as a navigation pointer. The argument does not depend on its content.\par\medskip
\ifdefined\hypertarget\hypertarget{sourcepointer-3}{}\fi
\noindent\textit{Source reference.} Original lemma lemma-realize-prescribed-residue-field-extension-etale: sources/stacks-more-morphisms/pinned-original-source.tex, source lines 10138--10145. Cited by the retained statement or proof but not part of the collected spans; supplied only as a navigation pointer. The argument does not depend on its content.\par\medskip
\ifdefined\hypertarget\hypertarget{sourcepointer-4}{}\fi
\noindent\textit{Source reference.} Original lemma 01TM: sources/stacks-more-morphisms/pinned-original-source.tex, source lines 12091--12091. Cited by the retained statement or proof but not part of the collected spans; supplied only as a navigation pointer. The argument does not depend on its content.\par\medskip
\ifdefined\hypertarget\hypertarget{sourcepointer-5}{}\fi
\noindent\textit{Source reference.} Original lemma 01TH: sources/stacks-more-morphisms/pinned-original-source.tex, source lines 12091--12091. Cited by the retained statement or proof but not part of the collected spans; supplied only as a navigation pointer. The argument does not depend on its content.\par\medskip
\ifdefined\hypertarget\hypertarget{sourcepointer-6}{}\fi
\noindent\textit{Source reference.} Original lemma morphisms-lemma-residue-field-quasi-finite: sources/stacks-more-morphisms/pinned-original-source.tex, source lines 12091--12091. Cited by the retained statement or proof but not part of the collected spans; supplied only as a navigation pointer. The argument does not depend on its content.\par\medskip
\ifdefined\hypertarget\hypertarget{sourcepointer-7}{}\fi
\noindent\textit{Source reference.} Original lemma schemes-lemma-points-fibre-product: sources/stacks-more-morphisms/pinned-original-source.tex, source lines 12091--12091. Cited by the retained statement or proof but not part of the collected spans; supplied only as a navigation pointer. The argument does not depend on its content.\par\medskip
\begin{lemma}
\label{lemma-etale-makes-quasi-finite-finite-multiple-points-var}
Let $f : X \to S$ be a morphism of schemes.
Let $x_1, \ldots, x_n \in X$ be points having the same image $s$ in $S$.
Assume that
\begin{enumerate}
\item $f$ is locally of finite type, and
\item $x_i \in X_s$ is isolated for $i = 1, \ldots, n$.
\end{enumerate}
Then there exist
\begin{enumerate}
\item[(a)] an \'etale neighbourhood $(U, u) \to (S, s)$,
\item[(b)] for each $i$ an integer $m_i$ and
open subschemes $V_{i, j} \subset X_U$, $j = 1, \ldots, m_i$
\end{enumerate}
such that we have
\begin{enumerate}
\item[(\romannumeral1)] each $V_{i, j} \to U$ is a finite morphism,
\item[(\romannumeral2)] there is a unique point $v_{i, j}$ of $V_{i, j}$
mapping to $u$ in $U$ with $\kappa(u) \subset \kappa(v_{i, j})$
finite purely inseparable,
\item[(\romannumeral4)] if $v_{i, j} = v_{i', j'}$, then $i = i'$ and
$j = j'$, and
\item[(\romannumeral3)] the points $v_{i, j}$ map to $x_i$ in $X$ and
no other points of $(X_U)_u$ map to $x_i$.
\end{enumerate}
\end{lemma}

\begin{proof}
This proof is a variant of the proof of
Algebra, Lemma \ref{algebra-lemma-etale-makes-quasi-finite-finite-variant}
in the language of schemes.
By Morphisms, Lemma \ref{morphisms-lemma-quasi-finite-at-point-characterize}
the morphism $f$ is quasi-finite at each of the points $x_i$.
Hence $\kappa(s) \subset \kappa(x_i)$ is finite for each $i$
(Morphisms, Lemma \ref{morphisms-lemma-residue-field-quasi-finite}).
For each $i$, let $\kappa(s) \subset L_i \subset \kappa(x_i)$
be the subfield such that $L_i/\kappa(s)$ is separable, and
$\kappa(x_i)/L_i$ is purely inseparable. Choose a finite Galois
extension $L/\kappa(s)$ such that there exist
$\kappa(s)$-embeddings $L_i \to L$ for $i = 1, \ldots, n$.
Choose an \'etale neighbourhood $(U, u) \to (S, s)$ such that
$L \cong \kappa(u)$ as $\kappa(s)$-extensions
(Lemma \ref{lemma-realize-prescribed-residue-field-extension-etale}).

\medskip\noindent
Let $y_{i, j}$, $j = 1, \ldots, m_i$ be the points of $X_U$
lying over $x_i \in X$ and $u \in U$. By
Schemes, Lemma \ref{schemes-lemma-points-fibre-product}
these points $y_{i, j}$ correspond exactly to the primes in the rings
$\kappa(u) \otimes_{\kappa(s)} \kappa(x_i)$. This also
explains why there are finitely many; in fact
$m_i = [L_i : \kappa(s)]$ but we do not need this.
By our choice of
$L$ (and elementary field theory)
we see that $\kappa(u) \subset \kappa(y_{i, j})$ is
finite purely inseparable for each pair $i, j$.
Also, by Morphisms, Lemma \ref{morphisms-lemma-base-change-quasi-finite}
for example, the morphism
$X_U \to U$ is quasi-finite at the points $y_{i, j}$ for
all $i, j$.

\medskip\noindent
Apply Lemma \ref{lemma-etale-makes-quasi-finite-finite-multiple-points}
to the morphism $X_U \to U$, the point $u \in U$
and the points $y_{i, j} \in (X_U)_u$. This gives an \'etale neighbourhood
$(U', u') \to (U, u)$ with $\kappa(u) = \kappa(u')$ and
opens $V_{i, j} \subset X_{U'}$ with the properties
(\romannumeral1), (\romannumeral2), and (\romannumeral3)
of that lemma. We claim that the \'etale neighbourhood
$(U', u') \to (S, s)$ and the opens $V_{i, j} \subset X_{U'}$
are a solution to the problem posed by the lemma.
We omit the verifications.
\end{proof}

\end{document}
